\usepackage{amsmath}
\usepackage{amssymb}
\usepackage{graphicx}
\usepackage{xcolor}
\usepackage{multicol}
\usepackage{caption}
\usepackage{eso-pic}
\setlength{\columnsep}{9mm}
\newcommand{\colsbegin}{\begin{multicols}{3}}
\newcommand{\colsend}{\end{multicols}}
% Level spreads (lib/format/registry.lua): the level opener sits alone on a right-hand
% (odd) page, the triptych on the facing left-hand page, the callout plate on the next
% right-hand page. \atlasopenright starts a right-hand page, leaving a blank left page
% when needed, even with openany for the rest of the book.
\newcommand{\atlasopenright}{\clearpage\ifodd\value{page}\else\thispagestyle{empty}\mbox{}\clearpage\fi}
\newenvironment{atlasopener}{\atlasopenright\thispagestyle{empty}\vspace*{\fill}\begin{center}\begin{minipage}{0.62\textwidth}}{\end{minipage}\end{center}\vspace*{\fill}\clearpage}
\newcommand{\atlasplatepage}[2]{\begin{center}\vspace*{\fill}\includegraphics[width=\textwidth,height=0.86\textheight,keepaspectratio]{#1}\par\vspace{4mm}{\small #2}\vspace*{\fill}\end{center}\clearpage}
% Large centred level title on the opener page (only inside the opener group).
% Level headings are written as \section* (unnumbered).
\makeatletter
\newcommand{\atlasopenertitle}{\renewcommand{\section}{\@startsection{section}{1}{0pt}{2mm}{8mm}{\centering\sffamily\bfseries\fontsize{34}{40}\selectfont}}}
\makeatother
